% GENERATED FILE. Edit canonical lessons, then run npm run generate:books.
% canonical-source-hash: fnv1a64:8eba38a4564920c8
% canonical-lessons: MW-C192-khatro, MW-C192-dhire, MW-C192-kharo, MW-C192-chhapno, MW-C192-sudharno

\chapter{Danger on the Road}
\label{ch:mw-danger-on-the-road}
\hlchaptermodality{192}

\begin{chapteropening}
\textbf{By the end of this chapter:} \emph{I can understand a warning by the road.}

Complete the last lesson of chapter 192: I can understand a warning by the road.
\end{chapteropening}

\section[\texorpdfstring{\mw{खतरो} (khatro)}{खतरो (khatro)}]{\mw{खतरो} (khatro) — danger}

\label{lesson:MW-C192-khatro}



\begin{quote}
Before the new one: say the Marwadi for careful, then the Marwadi for to obey.
\end{quote}

\subsection*{You'll want to know: \mw{खतरो}}
\textbf{\mw{खतरो}} — \emph{khatro} — \textquotedblleft{}danger\textquotedblright{}.

\begin{quote}
\mw{पुल} \mw{माथै} \mw{खतरो} \mw{है।}
\end{quote}

\begin{quote}
\emph{pul māthai khatro hai.}
\end{quote}

\begin{quote}
The bridge is dangerous.
\end{quote}

A masculine noun.

\begin{cousinweb}
\textbf{\mw{खतरो}} (\emph{khatro}) is usually traced, through Persian, to Arabic \emph{khaṭar}, \textquotedblleft{}danger\textquotedblright{}.
\end{cousinweb}

\begin{grammarlens}[title={What you've built}]
One more word for reading notices and messages.
\end{grammarlens}

\subsection*{Your turn}
\begin{itemize}
\raggedright
  \item \emph{Say it:} \emph{khatro}
  \item \emph{Say it:} \emph{khatro}, once more
  \item \emph{Say it:} \emph{pul māthai khatro hai.} — The bridge is dangerous.
  \item \emph{Recall:} say the Marwadi for careful, then the Marwadi for to obey, then say \emph{khatro} again
\end{itemize}

\begin{tcolorbox}[breakable,colback=teal!4,colframe=teal!35!black,title={Before you move on}]
What does \mw{खतरो} mean? (\textquotedblleft{}Danger\textquotedblright{}.) Say it once more.
\end{tcolorbox}
\section[\texorpdfstring{\mw{धीरे} (dhīre)}{धीरे (dhīre)}]{\mw{धीरे} (dhīre) — slowly}

\label{lesson:MW-C192-dhire}



\begin{quote}
Before the new one: say the Marwadi for to obey, then the Marwadi for danger.
\end{quote}

\subsection*{You'll want to know: \mw{धीरे}}
\textbf{\mw{धीरे}} — \emph{dhīre} — \textquotedblleft{}slowly\textquotedblright{}.

\begin{quote}
\mw{गळी} \mw{मांय} \mw{धीरे} \mw{चालो।}
\end{quote}

\begin{quote}
\emph{gaḷī mā̃y dhīre chālo.}
\end{quote}

\begin{quote}
Go slowly in the lane.
\end{quote}

\begin{cousinweb}
\textbf{\mw{धीरे}} (\emph{dhīre}) is inherited from Sanskrit \emph{dhīra}, \textquotedblleft{}steady, resolute\textquotedblright{}. CDIAL gives no Marwari form in this entry, but lists Sindhi \emph{dhīre}, \textquotedblleft{}quietly\textquotedblright{}, and Gujarati \emph{dhīrũ}, \textquotedblleft{}slow\textquotedblright{} (CDIAL 6817).
\end{cousinweb}

\begin{grammarlens}[title={What you've built}]
One more word for reading notices and messages.
\end{grammarlens}

\subsection*{Your turn}
\begin{itemize}
\raggedright
  \item \emph{Say it:} \emph{dhīre}
  \item \emph{Say it:} \emph{dhīre}, once more
  \item \emph{Say it:} \emph{gaḷī mā̃y dhīre chālo.} — Go slowly in the lane.
  \item \emph{Recall:} say the Marwadi for to obey, then the Marwadi for danger, then say \emph{dhīre} again
\end{itemize}

\begin{tcolorbox}[breakable,colback=teal!4,colframe=teal!35!black,title={Before you move on}]
What does \mw{धीरे} mean? (\textquotedblleft{}Slowly\textquotedblright{}.) Say it once more.
\end{tcolorbox}
\section[\texorpdfstring{\mw{खाड़ो} (khāṛo)}{खाड़ो (khāṛo)}]{\mw{खाड़ो} (khāṛo) — a pit; a pothole}

\label{lesson:MW-C192-kharo}



\begin{quote}
Before the new one: say the Marwadi for danger, then the Marwadi for slowly.
\end{quote}

\subsection*{You'll want to know: \mw{खाड़ो}}
\textbf{\mw{खाड़ो}} — \emph{khāṛo} — \textquotedblleft{}a pit; a pothole\textquotedblright{}.

\begin{quote}
\mw{सड़क} \mw{मांय} \mw{बड़ो} \mw{खाड़ो} \mw{है।}
\end{quote}

\begin{quote}
\emph{saṛak mā̃y baṛo khāṛo hai.}
\end{quote}

\begin{quote}
There's a big hole in the road.
\end{quote}

A masculine noun.

\begin{cousinweb}
\textbf{\mw{खाड़ो}} (\emph{khāṛo}) is inherited from a reconstructed \emph{*khāḍa}, \textquotedblleft{}a hollow\textquotedblright{}; CDIAL lists Marwari \emph{khāṛo}, \textquotedblleft{}hole\textquotedblright{} (CDIAL 3860).
\end{cousinweb}

\begin{grammarlens}[title={What you've built}]
One more word for reading notices and messages.
\end{grammarlens}

\subsection*{Your turn}
\begin{itemize}
\raggedright
  \item \emph{Say it:} \emph{khāṛo}
  \item \emph{Say it:} \emph{khāṛo}, once more
  \item \emph{Say it:} \emph{saṛak mā̃y baṛo khāṛo hai.} — There's a big hole in the road.
  \item \emph{Recall:} say the Marwadi for danger, then the Marwadi for slowly, then say \emph{khāṛo} again
\end{itemize}

\begin{tcolorbox}[breakable,colback=teal!4,colframe=teal!35!black,title={Before you move on}]
What does \mw{खाड़ो} mean? (\textquotedblleft{}A pit; a pothole\textquotedblright{}.) Say it once more.
\end{tcolorbox}
\section[\texorpdfstring{\mw{छापणो} (chhāpṇo)}{छापणो (chhāpṇo)}]{\mw{छापणो} (chhāpṇo) — to print}

\label{lesson:MW-C192-chhapno}



\begin{quote}
Before the new one: say the Marwadi for slowly, then the Marwadi for a pit.
\end{quote}

\subsection*{You'll want to know: \mw{छापणो}}
\textbf{\mw{छापणो}} — \emph{chhāpṇo} — \textquotedblleft{}to print\textquotedblright{}.

\begin{quote}
\mw{अखबार} \mw{रोज} \mw{सूचना} \mw{छापै} \mw{है।}
\end{quote}

\begin{quote}
\emph{akhbār roj sūchnā chhāpai hai.}
\end{quote}

\begin{quote}
The newspaper prints notices every day.
\end{quote}

Past (I …): \mw{म्हैं} \mw{सूचना} \mw{छापी}, \emph{mhaĩ sūchnā chhāpī}.

\begin{cousinweb}
\textbf{\mw{छापणो}} (\emph{chhāpṇo}) is inherited from a reconstructed \emph{*chapp}, \textquotedblleft{}press, cover, hide\textquotedblright{}; CDIAL lists Old Marwari \emph{chipāvaï}, \textquotedblleft{}hides\textquotedblright{}, and Hindi \emph{chāpnā}, \textquotedblleft{}to print\textquotedblright{}, under it (CDIAL 4994). To print is to press letters onto paper.
\end{cousinweb}

\begin{grammarlens}[title={What you've built}]
One more word for reading notices and messages.
\end{grammarlens}

\subsection*{Your turn}
\begin{itemize}
\raggedright
  \item \emph{Say it:} \emph{chhāpṇo}
  \item \emph{Say it:} \emph{chhāpṇo}, once more
  \item \emph{Say it:} \emph{akhbār roj sūchnā chhāpai hai.} — The newspaper prints notices every day.
  \item \emph{Recall:} say the Marwadi for slowly, then the Marwadi for a pit, then say \emph{chhāpṇo} again
\end{itemize}

\begin{tcolorbox}[breakable,colback=teal!4,colframe=teal!35!black,title={Before you move on}]
What does \mw{छापणो} mean? (\textquotedblleft{}To print\textquotedblright{}.) Say it once more.
\end{tcolorbox}
\section[\texorpdfstring{\mw{सुधारणो} (sudhārṇo)}{सुधारणो (sudhārṇo)}]{\mw{सुधारणो} (sudhārṇo) — to mend; to correct}

\label{lesson:MW-C192-sudharno}



\begin{quote}
Before the new one: say the Marwadi for a pit, then the Marwadi for to print.
\end{quote}

\subsection*{You'll want to know: \mw{सुधारणो}}
\textbf{\mw{सुधारणो}} — \emph{sudhārṇo} — \textquotedblleft{}to mend; to correct\textquotedblright{}.

\begin{quote}
\mw{मिस्त्री} \mw{सड़क} \mw{सुधारै} \mw{है।}
\end{quote}

\begin{quote}
\emph{mistrī saṛak sudhārai hai.}
\end{quote}

\begin{quote}
The mason is mending the road.
\end{quote}

Past (I …): \mw{म्हैं} \mw{साइकल} \mw{सुधारी}, \emph{mhaĩ sāikal sudhārī}.

\begin{cousinweb}
\textbf{\mw{सुधारणो}} (\emph{sudhārṇo}) is inherited from a reconstructed \emph{*śuddhakāra}, \textquotedblleft{}making clean\textquotedblright{}. CDIAL gives no Marwari form in this entry, but lists Hindi \emph{sudhārnā}, \textquotedblleft{}to mend\textquotedblright{} (CDIAL 12521).
\end{cousinweb}

\begin{grammarlens}[title={What you've built}]
One more word for reading notices and messages.
\end{grammarlens}

\subsection*{Your turn}
\begin{itemize}
\raggedright
  \item \emph{Say it:} \emph{sudhārṇo}
  \item \emph{Say it:} \emph{sudhārṇo}, once more
  \item \emph{Say it:} \emph{mistrī saṛak sudhārai hai.} — The mason is mending the road.
  \item \emph{Recall:} say the Marwadi for a pit, then the Marwadi for to print, then say \emph{sudhārṇo} again
\end{itemize}

\begin{tcolorbox}[breakable,colback=teal!4,colframe=teal!35!black,title={Before you move on}]
What does \mw{सुधारणो} mean? (\textquotedblleft{}To mend; to correct\textquotedblright{}.) Say it once more.
\end{tcolorbox}
